\documentclass[12pt]{article}
\usepackage[utf8]{inputenc}
\usepackage[T1]{fontenc}
\usepackage{hyperref}
\title{Tregami Language Resources for K-12 Curriculum Development}
\author{Compiled by Claude Sonnet 5 (claude-sonnet-5) for Andr\'as Kornai}
\date{}
\begin{document}
\maketitle

\section{Introduction}
Tregami (Trigami; ISO 639-3 \texttt{trm}) is a Nuristani language of the Indo-Iranian branch, spoken in the Tregam valley (villages including Gambir and Katar) in eastern Afghanistan. Sources place it in Kunar or Nuristan Province. Estimates are about 3,500 speakers (2011) against 1,000 in an older Ethnologue figure. It is closely related to Kalasha-ala (Waigali), and some sources question whether it is a separate language. The search found only general reference pages.

\section{Primary Educational Resources}
None were found in the search results.

\section{Linguistic and Reference Materials}
\begin{itemize}
\item Description of the Tregami on nuristan.info: \url{https://www.nuristan.info/Nuristani/Tregami/tregami.html}
\item Overview of Nuristani languages: \url{https://en.wikipedia.org/wiki/Nuristani_languages}
\item Language profile on Rosetta Panglossia: \url{https://rosettapanglossia.longnow.org/index.php?title=Trm}
\item Overview of the related Kalasha-ala (Waigali): \url{https://Www.Wikipedia.org/wiki/Waigali}
\end{itemize}

\section{Religious and Community Resources}
None were found in the search results.

\section{Digital and Audio Resources}
None were found in the search results. A Verbix page on Tregami appeared in the results, but I did not confirm what it holds.

\section{Curriculum Development Context}
No Tregami teaching materials, orthography or audio were found. The language is not known to be taught in schools. Curriculum work would start from documentation, with the related Kalasha-ala as a possible comparative reference.
\end{document}